\documentclass{article}
\usepackage[utf8]{inputenc}
\usepackage{array}
\newcolumntype{C}[1]{>{\centering\arraybackslash}m{#1}}
\usepackage{amsmath}
\usepackage{relsize}
\usepackage{graphicx}

\usepackage[noabbrev]{cleveref}

\crefformat{section}{\S#2#1#3}
\crefformat{subsection}{\S#2#1#3}
\crefformat{subsubsection}{\S#2#1#3}
\crefformat{figure}{#2Figura~#1#3}
\crefformat{footnote}{#2\footnotemark[#1]#3}
\crefformat{table}{#2Tabella~#1#3}
\crefformat{equation}{#2Equazione~#1#3}
\crefformat{chapter}{#2Capitolo~#1#3}

\begin{document}

\include{title}
\thispagestyle{empty}
\newpage

\section{Problema 1}
\paragraph{Traccia} Il palindromo è una sequenza di caratteri che, letta al contrario, rimane invariata. Ad
esempio, GAG e MADAM sono palindromi, ma ADAM no. Assumiamo che anche la stringa
vuota sia un palindromo. Da qualsiasi stringa non palindromica, si può sempre ottenere una
sotto-sequenza palindromica togliendo alcune lettere: ad esempio, data la stringa ADAM,
si può rimuovere la lettera M e ottenere ADA. Scrivere un programma per determinare la
lunghezza del palindromo più lungo che puoi ottenere da una stringa rimuovendo zero o
più caratteri. Il programma deve ricevere una stringa (di lunghezza inferiore a 1000
caratteri) e stampare la lunghezza del palindromo più lungo che si può ottenere rimuovendo
zero o più caratteri.

\paragraph{Soluzione}

La soluzione prevede l'utilizzo della programmazione dinamica. Invece di cercare tutti i possibili palindromi ricorsivamente, si osserva che una sotto-sequenza palindromica può essere trovata come una sotto-sequenza comune della stringa di partenza e l'inverso della stessa, ricordando che in un palindromo la prima e l'ultima lettera sono sempre uguali. In questo modo il problema si riduce a trovare la più lunga sotto-sequenza comune tra due stringhe. Per farlo, si è ricorso a una matrice di appoggio \textit{dp} che è stata popolata secondo le seguenti regole, scorrendo gli indici $i$ (riferito alla stringa di partenza) e $j$ (riferito a quella invertita):
\begin{itemize}
\item La prima riga e la prima colonna vengono inizializzate a 0;
\item Se i caratteri puntati dagli indici  (che indicano la prima e l'ultima lettera del possibile palindromo) non coincidono, allora non siamo in presenza di una sotto-sequenza comune e ci limitiamo a propagare il miglior risultato precedente facendo il massimo tra l'elemento superiore e quello a sinistra;
\item Se i caratteri puntati dagli indici coincidono, allora la sequenza continua a essere una sotto-sequenza comune e quindi viene incrementato il valore della lunghezza, aggiungendo uno rispetto al valore che si trova precedentemente sulla diagonale.
\end{itemize}

Definiamo quindi formalmente il problema della sotto-sequenza comune:

\[
    LCS(X_i, Y_j)=
\begin{cases}
    0,& \text{$i=0$ o $j=0$} \\
    max[LCS(X_{i}, Y_{j-1}),LCS(X_{i-1}, Y_{j})] & \text{$x_i \ne y_j$} \\
    LCS(X_{i-1}, Y_{j-1}), & \text{$x_i = y_j$}
\end{cases}
\]

La matrice verrà quindi riempita seguendo la definizione sopracitata come mostrato in \cref{fig:notebook}.

\begin{figure}[ht]
    \centering
    \includegraphics[width=13cm]{amogus-matrix.png}
    \caption{Popolamento matrice \textit{dp}}
    \label{fig:notebook}
\end{figure}

Memorizzando i valori nella matrice \textit{dp} evitiamo anche di dover ricalcolare le stesse sotto-sequenze che possono ricapitare.
Si noti come sono state aggiunte una riga e una colonna per soddisfare il caso $i=0$ o $j=0$.
La lunghezza massima della sotto-sequenza comune, che in questo caso coincide con quella palindromica, è data dal valore dell'ultima cella della matrice.

\paragraph{Complessità}
L'approccio brute force comporterebbe una complessità esponenziale, poiché bisognerebbe esaminare tutte le possibili combinazioni di sotto-sequenze esistenti. 
L'algoritmo proposto utilizza solo due cicli for innestati ottenendo quindi una complessità quadritica $O(n^2)$ dove $n$ è la lunghezza della stringa di partenza.
\end{document}
